\section*{Capítulo \ref{ch:g9:earth-resources} --- Os recursos da Terra}

\begin{solution}{exo:g9:earth-resources:1}
(a) \omterm{def:g9:earth-resources:aggregate}{Agregado}. (b) \omterm{def:g9:earth-resources:reserve}{Reserva}. (c) \omterm{def:g9:earth-resources:ore-deposit}{Jazida}. (d) \omterm{def:g9:earth-resources:ore-deposit}{Fator de concentração}.
\end{solution}

\begin{solution}{exo:g9:earth-resources:2}
(a) Nos \omterm{def:g3:stones-pebbles-sand:pebble}{cascalhos} dos \omterm{def:g1:landscapes-around-us:water}{rios}. (b) Nas fendas, pela água quente. (c) No
\omterm{def:g7:volcanoes:magma}{magma}. (d) É a primeira etapa de um depósito de \omterm{def:g3:stones-pebbles-sand:pebble}{cascalho} de \omterm{def:g1:landscapes-around-us:water}{rio}: o ouro
liberado é depois levado pelo \omterm{def:g1:landscapes-around-us:water}{rio}.
\end{solution}

\begin{solution}{exo:g9:earth-resources:3}
(b), (d), (a), (e), (c).
\end{solution}

\begin{solution}{exo:g9:earth-resources:4}
(a) Não renovável: leva milhões de anos para se formar. (b) Renovável: a
\omterm{def:g4:clouds-rain-snow:precipitation}{chuva} repõe o que é retirado. (c) Não renovável: uma \omterm{def:g9:earth-resources:ore-deposit}{jazida} se forma em
tempos muito longos. (d) Não renovável nesse ritmo: o \omterm{def:g6:water-underground:water-table}{lençol freático}
baixa e os \omterm{def:g6:water-underground:spring}{poços} acabam secando.
\end{solution}

\begin{solution}{exo:g9:earth-resources:5}
(a) $4 \div 0.007 \approx 571$, cerca de $10^3$ (entre $10^2$ e $10^3$,
mais perto de $10^3$). (b) $0.5 \div 0.0002 = 2500$, cerca de $10^3$.
(c) $20
\div 5.6 \approx 3.6$, cerca de 4: entre $10^0$ e $10^1$.
(d) O ferro já é muito abundante na \omterm{def:g6:inside-the-earth:layers}{crosta}: um minério só algumas vezes
mais rico que a \omterm{def:g3:stones-pebbles-sand:rock}{rocha} comum basta.
\end{solution}

\begin{solution}{exo:g9:earth-resources:6}
(a) $19.3 \div 2.5 \approx 7.7$ vezes. (b) Ouro: $19.3 \times 0.01 =
\qty{0.193}{g}$; \omterm{def:g3:stones-pebbles-sand:rock}{rocha}: $2.5 \times 0.01 = \qty{0.025}{g}$. (c) Quando
todo o leito se move, os grãos pesados de ouro afundam entre a \omterm{def:g3:stones-pebbles-sand:sand}{areia} e os
\omterm{def:g3:stones-pebbles-sand:pebble}{seixos} mais leves, até a \omterm{def:g3:stones-pebbles-sand:rock}{rocha} de base sólida, onde param.
\end{solution}

\begin{solution}{exo:g9:earth-resources:7}
(a) O ouro, o estanho e o cromo. (b) Entre $10^1$ e $10^2$, perto de
$10^2$. (c) $3000 \div 3.7 \approx 800$ vezes. (d) O cromo (ou o
estanho, ou o ouro): quanto maior o \omterm{def:g9:earth-resources:ore-deposit}{fator de concentração}, mais raros os
lugares onde a natureza concentrou o metal o bastante.
\end{solution}

\begin{solution}{exo:g9:earth-resources:8}
(a) $18 \times 365 = \qty{6570}{kg}$, cerca de \qty{6.6}{t}. (b) $5 \times
10^{10}$ toneladas. (c) $25000 \div 18 \approx 1389$ dias, quase quatro
anos. (d) A \omterm{def:g3:stones-pebbles-sand:sand}{areia} tirada do \omterm{def:g1:landscapes-around-us:water}{rio} não desce mais até o \omterm{def:g1:landscapes-around-us:water}{litoral}; as ondas
continuam levando \omterm{def:g3:stones-pebbles-sand:sand}{areia} da praia, mas chega menos \omterm{def:g3:stones-pebbles-sand:sand}{areia} para repô-la.
\end{solution}

\begin{solution}{exo:g9:earth-resources:9}
(a) De minúsculas plantas e minúsculos animais do \omterm{def:g2:water-on-earth:ocean}{mar}, como as
diatomáceas, enterrados na \omterm{def:g3:stones-pebbles-sand:sand}{lama} do fundo do \omterm{def:g2:water-on-earth:ocean}{mar}. (b) Entre cerca de 60 e
\qty{120}{\degreeCelsius}. (c) Eles são mais leves que a água, que enche
os \omterm{def:g6:water-underground:porosity}{poros} das \omterm{def:g3:stones-pebbles-sand:rock}{rochas}. (d) Uma camada \omterm{def:g6:water-underground:permeability}{impermeável}, como o folhelho; eles se
acumulam embaixo dela, muitas vezes no alto de uma dobra arqueada para
cima, o gás acima do petróleo.
\end{solution}

\begin{solution}{exo:g9:earth-resources:10}
(a) $2.9 \times 10^7$ toneladas. (b) $2.9 \times 10^7 \div 100 = 2.9 \times
10^5$ toneladas por ano, mais ou menos. (c) O minério é 100 vezes mais
pesado que o cobre que contém: $2.9 \times 10^7 \times 100 = 2.9 \times 10^9$ toneladas.
(d) Cada tonelada de minério dá só \qty{10}{kg} de cobre, então bilhões
de toneladas de \omterm{def:g3:stones-pebbles-sand:rock}{rocha} têm de ser escavadas.
\end{solution}

\begin{solution}{exo:g9:earth-resources:11}
(a) A 23~milhões de toneladas por ano, ela teria durado cerca de $280 \div 23
\approx 12$ anos, até o início dos anos 1980; mas foram lavradas
697~milhões de toneladas, muito mais do que essa \omterm{def:g9:earth-resources:reserve}{reserva}. (b) $980 \div 280 = 3.5$ vezes.
(c) $980 \div
23 \approx 43$ anos. (d) Por exemplo: novas \omterm{def:g9:earth-resources:ore-deposit}{jazidas} são descobertas;
técnicas melhores ou preços mais altos fazem \omterm{def:g9:earth-resources:ore-deposit}{jazidas} mais pobres valerem
a pena.
\end{solution}

\begin{solution}{exo:g9:earth-resources:12}
(a) $5 \times 10^1$ anos. (b) $2 \times 10^3 \div 50 = 40$ vezes. (c) $3
\times 10^8 \div (5 \times 10^1) = 6 \times 10^6$ vezes. (d) O carvão
levou milhões de vezes mais tempo para se formar do que levamos para
usá-lo: a natureza não o renova na nossa escala de tempo.
\end{solution}

\begin{solution}{pb:g9:earth-resources:1}
\textbf{1.} $100 \times 1000 = \num{100000}$~kg $= 1 \times 10^5$~kg.

\textbf{2.} $2000 \times 300 = \num{600000}$~t $= 6 \times 10^5$~t.

\textbf{3.} $6 \times 10^5 \div 10 = 6 \times 10^4$~t por ano.

\textbf{4.} $0.40 \times 6 \times 10^4 = 2.4 \times 10^4$~t de \omterm{def:g3:stones-pebbles-sand:sand}{areia} por
ano.

\textbf{5.} $1200 \times 500 \times 5 = 3 \times 10^6$~m$^3$.

\textbf{6.} Dois terços dela: $2 \times 10^6$~m$^3$.

\textbf{7.} $2 \times 10^6 \times 1.6 = 3.2 \times 10^6$~t.

\textbf{8.} $\num{24000} + \num{136000} = \num{160000}$~t $= 1.6 \times
10^5$~t por ano.

\textbf{9.} $\num{160000} \div 2000 = 80$ vezes mais.

\textbf{10.} Não: a cava é esvaziada 80 vezes mais depressa do que o \omterm{def:g1:landscapes-around-us:water}{rio}
renova a \omterm{def:g3:stones-pebbles-sand:sand}{areia}. Na escala da cidade, essa \omterm{def:g3:stones-pebbles-sand:sand}{areia} é um recurso finito.

\textbf{11.} $\num{24000} \div 4 = \num{6000}$~t de \omterm{def:g3:stones-pebbles-sand:sand}{areia} por ano.

\textbf{12.} $3.2 \times 10^6 \div (1.6 \times 10^5) = 20$: \textbf{20 anos}.
\end{solution}
